% ============================================================
%  Chương I – Giới thiệu
%  ~300 từ · 1 trang
% ============================================================
\chapter{Giới thiệu}
\label{chap:intro}

% Ý chính: đặt câu hỏi trung tâm — mật mã bất đối xứng giải quyết được
% key exchange nhưng chưa giải quyết identity binding. PKI là câu trả lời.

% ----------------------------------------------------------
\section{Vấn đề đặt ra}
\label{sec:1.1}

% Ý chính:
%   - Mật mã khóa công khai cho phép trao đổi bí mật và ký số mà không cần
%     chia sẻ khóa trước, nhưng chưa trả lời: "Khóa công khai này có thực sự
%     thuộc về đúng entity đó không?"
%   - Kẻ tấn công MITM có thể thay thế khóa công khai trên kênh truyền
%     mà hai đầu không hay biết
%   - Cần cơ chế bên thứ ba tin cậy để ràng buộc (identity ↔ public key)

% ----------------------------------------------------------
\section{Phạm vi và bố cục báo cáo}
\label{sec:1.2}

% Ý chính: định hướng đọc — mỗi chương trả lời một câu hỏi cụ thể
%   - Chương~\ref{chap:crypto}   : tại sao cần PKI? (nền tảng lý thuyết)
%   - Chương~\ref{chap:pki-arch}: PKI gồm những thành phần nào?
%   - Chương~\ref{chap:x509}    : chứng chỉ số được định nghĩa như thế nào?
%   - Chương~\ref{chap:deploy}  : PKI được dùng ở đâu trong thực tế?
%   - Chương~\ref{chap:opensource}: có thể tự triển khai PKI bằng gì?
%   - Chương~\ref{chap:challenges}: vấn đề còn tồn đọng và hướng tương lai
